%%% lorem.tex --- 
%% 
%% Filename: lorem.tex
%% Description: 
%% Author: Ola Leifler
%% Maintainer: 
%% Created: Wed Nov 10 09:59:23 2010 (CET)
%% Version: $Id$
%% Version: 
%% Last-Updated: Tue Oct  4 11:58:17 2016 (+0200)
%%           By: Ola Leifler
%%     Update #: 7
%% URL: 
%% Keywords: 
%% Compatibility: 
%% 
%%%%%%%%%%%%%%%%%%%%%%%%%%%%%%%%%%%%%%%%%%%%%%%%%%%%%%%%%%%%%%%%%%%%%%
%% 
%%% Commentary: 
%% 
%% 
%% 
%%%%%%%%%%%%%%%%%%%%%%%%%%%%%%%%%%%%%%%%%%%%%%%%%%%%%%%%%%%%%%%%%%%%%%
%% 
%%% Change log:
%% 
%% 
%% RCS $Log$
%%%%%%%%%%%%%%%%%%%%%%%%%%%%%%%%%%%%%%%%%%%%%%%%%%%%%%%%%%%%%%%%%%%%%%
%% 
%%% Code:

\chapter{Theory}
\label{cha:theory}

Applications of AI led to improvements in different areas in healthcare \todo{enter examples}. The quality of these solutions depends on the data 
that they are trained on. For most models it has to have a high volume and quality. Due to different issues, such datasets are seldom available: 
privacy issues and limited sharing between hospitals and research institutions make it hard to gather medical data in the first place. Then, 
annotating the data is expensive, as trained medical personal is required to assure a high quality of labels \cite{ibrahim_generative_2025}.

An approach to create bigger and high quality datasets is to use generative models to create synthetic data that is highly similar to real data\todo{include related work that showed this}.    

This chapter explains the foundations of generative models and different architectures. It also covers the theory behind image classification models
that are used to evaluate the capability to use the synthetic data to augment real data for a downstream task. 

\section{Generative Models}

The general goal of generative models is to approximate an unknown distribution $p_{data}(\boldsymbol{x})$ of the training data $\boldsymbol{x}$ 
with a learned distribution $p_{model}(\boldsymbol{x})$ \cite{goodfellow_generative_2020}. $p_{model}$ can then be used to draw new samples, that 
could also come from $p_{data}$. There are different ways to train a model to learn this distribution. The most popular approaches are Generative 
Adversarial Networks (GANs) and Diffusion Models (DMs).

Different types of input can be given to the models to create a synthetic image. The basic case is to use noise as a starting point. The noise 
vector $\boldsymbol{z}$ follows a distribution $p(\boldsymbol{z})$, usually a high-dimensional Gaussian. These models are called noise-to-image 
models. Image-to-image models take, as the name indicates, another image as input. It is used to translate between the domain of the input image 
and the output image and finds application in e.g. image inpainting, image colorization or to increase the resolution of an image. Popular are also 
text-to-image models that can create images from a text prompt.


noise-to-image, image-to-image, text-to-image; medical context

\subsection{Generative Adversarial Networks}

For a long time, GANs dominated the field of synthetic image generation. They were introduced by Goodfellow et al. \cite{goodfellow_generative_2020}
in 2014 and are the foundation for many adjustments and improvements until today. 

The basic idea is to have two adversarial networks that try to outperform each other. On the one side stands the generator 
$G(\boldsymbol{z}; \boldsymbol{\Theta}^{G})$. $z$ is the input noise vector and the generator tries to learn how transform it into a realistic image. 
$\boldsymbol{\Theta}^{G}$ is a set of parameters that are learned by the generator to improve its 
ability to fool its counter part, which is the discriminator $D(\boldsymbol{z}; \boldsymbol{\Theta}^{G})$. $\boldsymbol{\Theta}^{D}$. 
The discriminator tries to classify input correctly as real or synthetic images. It also has a set of parameters $\boldsymbol{\Theta}^{D}$ that are 
learned during the training process.

An adversarial loss function is defined (see formula \ref{eq:adv_loss}), that the generator tries to minimize and the discriminator tries to maximize. 

\begin{equation}
    \min_G \max_D 
    [
    \mathbb{E}_{x \sim p_{\text{data}}(x)}[\log D(x)] +
    \mathbb{E}_{z \sim p_z(z)}[\log(1 - D(G(z)))]
    ]
    \label{eq:adv_loss}
\end{equation}



\subsection{Diffusion Models}

from first idea to CycleDiff

\subsection{Metrics}

FID/IS

Radiometrics + statistical test



\section{Image Classification}

few general words about history of image classification, in general and medical context

\subsection{Convolutional Neural Networks}

\subsection{Metrics}


%%%%%%%%%%%%%%%%%%%%%%%%%%%%%%%%%%%%%%%%%%%%%%%%%%%%%%%%%%%%%%%%%%%%%%
%%% lorem.tex ends here

%%% Local Variables: 
%%% mode: latex
%%% TeX-master: "demothesis"
%%% End: 
